% !TEX root = ../paper.tex

\section{Protocol Definitions and Reproducibility Details}
\label{sec:protocol_details}

\subsection{Probe Frames and Relative Quality Signals}

Each scene uses up to 12 QA probes sampled evenly from its full manifest pair list.
Images are evaluated at quarter scale with the longest
dimension capped at 640 pixels. The scene homography is converted to these evaluation coordinates.
For each probe, both the transformed image and the size-matched reference are scored within the
same warped infrared validity mask. The reference resizes infrared edges and gradients to the RGB
evaluation grid. The resize-only mapping is denoted by $H_0$.

Let $E_i(H)$ denote edge-overlap F1 with a one-pixel dilation tolerance, and $G_i(H)$ denote
normalized cross-correlation between gradient magnitudes inside the validity mask. Define
\begin{equation}
\delta E_i = E_i(\bar H_S)-E_i(H_0), \qquad
\delta G_i = G_i(\bar H_S)-G_i(H_0).
\end{equation}
The scene gains $\Delta E_S$ and $\Delta G_S$ are probe means. A severe relative outlier has both
$\delta E_i<-0.01$ and $\delta G_i<-0.01$; $O_S$ is its fraction among probes. Let $p_{\lor}$ be
the fraction with either gain positive and $p_G$ the fraction with positive gradient gain.
These image-domain diagnostics measure agreement relative to the size-matched reference.

\subsection{Support and Canonical Verification}
\label{sec:canonical_details}

A frame enters the accepted pool when homography fitting succeeds with at least 80 candidate
matches, inlier ratio at least 0.30, mean inlier reprojection error at most 4 pixels, and coverage
at least 0.25. For a scene with $N$ valid pairs, the frozen minimum accepted support is
\begin{equation}
N_{\min}=\max\!\left(4,\min\!\left(N,\max\!\left(8,
\min\!\left(20,\operatorname{round}(0.15N)\right)\right)\right)\right).
\end{equation}
This yields thresholds of 8 for S01--S04, 15 for S07, and 20 for the remaining scenes.
Accepted support $A_S$ is the number of accepted frames divided by the number attempted.
The robust-rejection ratio $J_S$ is the fraction removed from this accepted pool during consensus.
Dispersion $D_S$ is the median, across accepted hypotheses, of mean displacement from the consensus
on a $5\times5$ infrared-image grid, measured in RGB-image pixels.

The primary path $L_S$ requires $\Delta E_S>0$, $\Delta G_S>0$, $p_{\lor}\ge0.60$, and zero
severe relative outliers. The alternative path $W_S$ applies to infrared input with size-based
initialization and requires all of the following:
\begin{equation}
\begin{aligned}
A_S&\ge0.80, & \bar r_S&\ge0.30, & \bar e_S&\le4\text{ px}, & \bar c_S&\ge0.25, \\
D_S&\le25\text{ px}, & J_S&\le0.25, & O_S&\le0.10, & n_{\mathrm{out}}&\le1, \\
\Delta E_S&\ge-0.05, & \Delta G_S&>0, & p_G&\ge0.60. &&
\end{aligned}
\end{equation}
Here $\bar r_S,\bar e_S,\bar c_S$ are accepted-frame means. Path $L_S$ gives a pass;
path $W_S$ gives a pass with warning when $L_S$ is not met. Both contribute to canonical retention
once $N_S\ge N_{\min}$. Adaptive expansion stops when this retention condition and the
$D_S\le25$, $J_S\le0.25$ stability checks hold. After exhausting the schedule, the final
canonical decision still follows Eq.~\ref{eq:canonical_gate}; the stopping checks are not an
additional requirement on the primary path.

\subsection{Normalized Reliability Score}
\label{sec:score_normalization}

Let $C(x)=\min(1,\max(0,x))$. The components used in Eq.~\ref{eq:reliability_score} are
\begin{equation}
\begin{aligned}
\hat a_S &= C(A_S), & \hat o_S &= 1-C(O_S/0.10), \\
\hat j_S &= 1-C(J_S/0.25), & \hat e_S &= C((\Delta E_S+0.05)/0.20), \\
\hat g_S &= C(\Delta G_S/0.15), &&
\end{aligned}
\end{equation}
\begin{equation}
\hat z_S=0.5C(\bar r_S/0.50)+0.3C(\bar c_S)
       +0.2\left[1-C\!\left((\bar e_S-1.5)/1.0\right)\right].
\end{equation}
Fixed constants define a diagnostic ranking score in $[0,100]$, not a calibrated probability.
Score-ranked coverage points select
the first $K$ scenes in this ordering. The canonical point is evaluated on its own retained set.

\subsection{Metric Aggregation}
\label{sec:metric_aggregation}

\paragraph{Pairwise statistics.}
Inlier ratio is the number of fitted inliers divided by the candidate-match count. Coverage is
the fraction of the 16 infrared-image grid cells occupied by candidate matches, before robust
fitting. Reprojection error is the mean Euclidean inlier residual in RGB-image pixels for one
pair; the main table reports the median of these pair-level means. Each summary uses the recorded
finite, non-null values of that metric. Recorded zero inlier ratios and coverages are included,
whereas missing reprojection values are excluded. Runtime is the mean recorded time per pair
under the stated implementation and deployment profile.

\paragraph{Scene and condition statistics.}
Scene coverage is the fraction of scenes retained. Pair coverage weights each retained scene by
its valid manifest pair count. Accepted/attempted ratios use the summed frame counts within the
indicated scene set. Accordingly, retained-set support is $1491/2304=64.7\%$, whereas all-scene
support is $2473/3847=64.3\%$. Condition retention is scene-macro; accepted support is frame-micro.

\paragraph{Resampling and perturbations.}
For leave-frame-out analysis, grid displacement is averaged within each fold, then across the
five folds of each scene. The reported 1.39 pixels is the median of these 15 scene means;
their mean is 2.26 pixels. Perturbation curves average the joint edge/gradient score change
over 45 frames at each severity and mode. A degradation rate counts frames with a negative joint
score change at that setting. These tests assess consensus stability and quality-signal response,
respectively.
